% Part: proof-theory
% Chapter: sequent-calculus
% Section: introduction

\documentclass[../../../include/open-logic-section]{subfiles}

\begin{document}

\olfileid[tr]{pt}{seq}{int}

\olsection{Giriş}

Sekant hesapları, Gerhard Gentzen'in 1930'larda tanıttığı !!{proof}
sistemleridir. Bunlar daha çok şeyleri !!{proving}{}k için pratik bir
yöntem olmaktan ziyade, Gentzen'in !!{proof}s{}ın olanaklı yapısı ve özellikleri
hakkında belirli sonuçlar ispatlayabilmesi için tasarlanmıştı. Bu tür sorular
tam da ispat kuramının ele aldığı sorulardır.

Sekant hesapları, adından da anlaşılacağı üzere !!{formula}s üzerinde değil,
\emph{sekantlar} üzerinde işler. Bir sekant, iki !!{formula}s dizisinin ya da
çoklukümesinin oluşturduğu çifttir. Gentzen'in özgün sistemleri ($\Log{LK}$ ve
$\Log{LJ}$) diziler kullanıyordu; günümüzde ispat kuramcıları çoklukümeleri
yeğler. Çokluküme bir kümeye benzer; ancak bir !!{element}{}ı birden çok kez
içerebilir. Yine de kümelerde olduğu gibi !!{element}s{}ın sırası önemli
değildir. Dolayısıyla $\{!A,!A,!B\}$ ile $\{!A,!B,!A\}$ ifadeleri aynı
çoklukümeyi anlatır. Buna karşılık o çokluküme $\{!A,!B\}$ ve
$\{!A,!A,!B,!B\}$ çoklukümelerinden farklıdır; çünkü bunlar öncekinden farklı
sayıda $!A$ ve $!B$ içerir.

İspat kuramında !!{formula}s çoklukümeleri geleneksel olarak büyük Yunan
harfleriyle yazılır. Dolayısıyla $\Gamma$, $\Delta$, $\Pi$, $\Lambda$ ya da
$\Theta$ yazdığımızda bunların çalıştığımız mantığın !!{formula}s{}inden oluşan
çoklukümeler olmasını amaçlarız. İspat kuramcıları yine geleneksel olarak
$\tuple{\Gamma,\Delta}$ yazmak yerine iki bileşeni bir sekant simgesiyle
ayırır; biz~$\Sequent$ kullanıyoruz.

\begin{defn}[Sekant]
Bir \emph{sekant},
\[
\Gamma \Sequent \Delta
\]
biçimindeki bir ifadedir; burada $\Gamma$ ile $\Delta$, sonlu (boş da
olabilen) !!{formula}s çoklukümeleridir. $\Gamma$'ya \emph{önbileşen},
$\Delta$'ya \emph{artbileşen} denir.
\end{defn}

$!G$, $\Gamma$ içindeki !!{formula}s{}in tümel evetlemesi (boş tümel evetleme
$\ltrue$ alınır); $!D$ ise $\Delta$ içindeki !!{formula}s{}in ayrışımı (boş
ayrışım $\lfalse$ alınır) olsun. O zaman $\Gamma \Sequent \Delta$ sekantı,
(bir $\Struct{M}$ yapısında) ancak ve ancak $!G \lif !D$ doğruysa doğrudur.
Klasik mantıkta bu, $\Gamma$ içindeki !!{formula}s{}den en az birinin yanlış ya
da $\Delta$ içindeki !!{formula}s{}den en az birinin doğru olmasıyla aynıdır.

$\Gamma$ bir !!{formula}s çoklukümesiyse $\Gamma,!A$ (ya da $!A,\Gamma$),
$\Gamma$'ya $!A$ eklenmesiyle elde edilen çoklukümeyi gösterir. $\Delta$ da
bir !!{formula}s çoklukümesiyse $\Gamma,\Delta$ iki çoklukümenin
çoklukümesel birleşimidir: $\Gamma$, $!A$'yı $n$ \emph{çokluğuyla} (yani
$n$ kopya $!A$) ve $\Delta$, $!A$'yı $m$ çokluğuyla içeriyorsa
$\Gamma,\Delta$, $!A$'yı $n+m$ çokluğuyla içerir. Sekantın bir tarafındaki
çokluküme boşsa genellikle o tarafı boş bırakırız. Örneğin $\Sequent !A$,
önbileşeni boş ve artbileşeni $!A$'yı $1$ çokluğuyla içeren sekanttır.

Sekant hesabındaki !!^a{proof}, sekantlardan oluşan bir ağaçtır. En üstteki
(ya da başlangıç) sekantlar ele alınan sistemin aksiyomlarıdır; bir sekant
başka bir ya da iki sekantın altında bulunuyorsa, sistemin bir çıkarım
kuralıyla onlardan doğru biçimde çıkmalıdır. Önce klasik mantık için iki ayrı
sekant sistemini, $\Log{G1c}$ ile $\Log{G3c}$'yi ele alacağız. Aksiyomları ve
kuralları \cref{pt:seq:int:tab:G1c,pt:seq:int:tab:G3c} içinde verilmiştir.
\oliflabeldef{fol:seq::chap}{ (Gentzen'in özgün $\Log{LK}$ sistemi
\olref[fol][seq][]{chap} içinde açıklanır.)}{}

\subfile{rules-G1c}
\subfile{rules-G3c}

Sistemler ince biçimde farklıdır ve bu farklar değişik ispat kuramsal
özelliklere sahip olmalarına yol açar. $\Log{G1c}$, başka hiçbir !!{formula}
içermesine izin verilmeyen $!A \Sequent !A$ aksiyomlarına sahiptir.
$\Log{G3c}$ ise $\Gamma$ ile $\Delta$ içinde gelişigüzel ``yan'' !!{formula}s
bulunan $!C,\Gamma \Sequent \Delta,!C$ biçiminde aksiyomlara sahiptir; fakat
iki tarafta da görünen !!{formula}~$!C$ \emph{atomik} olmalıdır. $\Log{G1c}$,
$\LeftR{\land}$ ile $\RightR{\lor}$ kurallarının ikişer sürümünü;
$\Log{G3c}$ ise birer sürümünü içerir. Ayrıca $\Log{G1c}$, ``büzülme''
(\LeftR{\Contraction}, \RightR{\Contraction}) ve ``zayıflatma''
(\LeftR{\Weakening}, \RightR{\Weakening}) denen ek ``yapısal kurallara''
sahiptir.

$\Log{G1c}$ ile $\Log{G3c}$ klasik mantık için sağlam ve tamdır. Başka bir
deyişle bu sistemlerde !!{provable} olan her sekant her !!{structure} içinde
doğrudur ve her !!{structure} içinde doğru olan her sekantın !!a{proof}{}ı
vardır. Dolayısıyla bir sekantın !!a{proof}{}ının bulunup bulunmadığı sorusu,
mantığın başka yöntemleriyle de (örneğin model kuramıyla) yanıtlanabilir. İki
sistem de tam olarak her !!{structure} içinde doğru olan sekantları !!{prove}{}dığından
özellikle aynı sekantları ispatlar.

Ancak ispat kuramı yalnızca bir şeyin ispatlanabilir \emph{olup olmadığıyla}
değil, \emph{nasıl} ispatlandığıyla da ilgilenir. İspat kuramcıları ``Belirli
bir kuraldan her zaman kaçınılabilir mi?'', ``Bu kural, yeni sekantları
ispatlanabilir kılmadan sisteme eklenebilir mi?'' ya da ``Bir sistemdeki
!!{proof}s başka bir sistemdeki !!{proof}s{}a dönüştürülebilir mi?'' gibi
soruları ele alır. Böyle bir sorunun yanıtı ``evet'' ise olgunun ispatı çoğu
kez !!{proof}s{}ın karmaşıklığı üzerinde, denk biçimde !!{proof}s{}ın yapısı
üzerinde tümevarımla verilir. Bu yalnızca olgunun ispatını (örneğin
$\Log{G1c}$ bir sekantı ispatlıyorsa $\Log{G3c}$'nin de aynı sekantı
ispatladığını) değil, bir yordamı da (örneğin $\Log{G1c}$ içindeki !!{proof}s{}ı
$\Log{G3c}$ içindeki !!{proof}s{}a dönüştürmeyi) verir.

İspat kuramının merkezî bir sonucu \emph{kesme-giderimi teoremi}dir. Bu
teorem, sekant sistemlerine hangi sekantların !!{provable} olduğunu
değiştirmeden
\[
\Axiom$\Gamma \fCenter \Delta, !A$
\Axiom$!A, \Pi \fCenter \Lambda$
\RightLabel{\Cut}
\BinaryInf$\Gamma, \Pi \fCenter \Delta, \Lambda$
\DisplayProof
\]
kesme kuralını ekleyebileceğimizi gösterir. Dahası ispat, $\Log{G1c}$'nin
(ya da $\Log{G3c}$'nin) kurallarıyla birlikte \Cut{} kullanan her !!{proof}{}ı,
yalnızca $\Log{G1c}$'nin (ya da $\Log{G3c}$'nin) kurallarını kullanan bir
ispata dönüştüren yöntem verir. Başka bir deyişle yordam, \Cut{} kuralını onu
kullanan !!{proof}s{}dan ``giderir''; adı buradan gelir.

\end{document}
